\documentclass[11pt]{article}
\usepackage[utf8]{inputenc}
\usepackage{geometry}
\geometry{a4paper, margin=1in}
\usepackage{titlesec}
\usepackage{enumitem}
\usepackage{xcolor}

\title{\textbf{Strategic Diagnostic Snapshot™}}
\author{RankPilot Intelligence}
\date{\today}

\begin{document}

\maketitle

\section*{I. Executive Summary}
\textbf{Firm:} González de Araujo Consultores \\
\textbf{Practice Area:} Data Protection \\
\textbf{Detected Practice Model:} Standard

\section*{II. Structural Positioning}
\textbf{Current Tier Assessment:} Unranked / New Entry \\
\textbf{Strategic Confidence:} 100.0\% \\
\textbf{Advantage:} Standard Market Offerings

\section*{III. Portfolio Analysis (Key Matters)}
\begin{itemize}
\item \textbf{Grupo Hermes Data Protection Framework} (Client: Grupo Hermes | Value: N/A)\\ Acting significance of the Grupo Hermes, a Mexican diversified infrastructure and services conglomerate, Data Protection Framework lies in its documented conversion of data protection and governance requirements into operational processes with a quantified result: 100% compliance in handling ARCO requests and the avoidance of regulatory sanctions. Rather than being limited to the preparation of isolated documents, the mandate concerned the design and implementation of a thorough data protection and governance framework that linked the client’s personal-data handling activities to defined governance measures.

Grupo Hermes, developing tailored privacy notices, implementing internal policies, and establishing procedures for ARCO rights management, engaged the firm in an exercise requiring detailed identification of how personal data moved through the organisation and how ARCO rights would be managed in practice. The firm played a central role in mapping personal data flows, developing tailored privacy notices, implementing internal policies, and establishing procedures for ARCO rights management. This structured approach brought together the mapping of personal data flows, tailored privacy notices, internal policies and ARCO rights-management procedures as connected components of the data protection and governance framework.

Arturo González de Araujo and Ernesto Zavala led the engagement. Their work supported the design and implementation of the thorough framework and the practical establishment of procedures for ARCO rights management. The outcome recorded for Grupo Hermes was 100% compliance in handling ARCO requests and avoiding regulatory sanctions, providing concrete evidence of the framework’s operational application.
\item \textbf{Mega Direct Data Protection Advisory} (Client: MEGA DIRECT, S.A. de C.V. | Value: N/A)\\ Advising for MEGA DIRECT, S.A. de C.V., a customer experience, call center, and marketing services provider, over 16 years, the team’s work demonstrates sustained counsel on the legal treatment of databases as business-critical assets, rather than a discrete or one-off compliance exercise. The duration of the advisory relationship is a material indicator of the client’s continued need for advice on complex data protection and database-related matters.

MEGA DIRECT, S.A. de C.V. instructed the firm in connection with legal proceedings concerning data acquisition, processing, storage, and protection. Those proceedings required the team to address the full lifecycle of the relevant data: how it was acquired, processed, stored and protected, and the legal consequences arising from each of those stages. The mandate therefore placed database-related legal issues directly within the context of proceedings affecting the client’s ongoing operations.

Arturo González de Araujo, Ernesto Zavala and Damián Ortega led the representation. Their involvement across the 16-year advisory period reflects a sustained engagement with complex data protection and database-related matters, including representation in the legal proceedings concerning data acquisition, processing, storage, and protection. The work ensured business continuity in a highly regulated environment, preserving the operational relevance of the client’s data assets while those issues were addressed through the relevant legal proceedings.
\item \textbf{Biocodex Data Protection Framework} (Client: Biocodex | Value: N/A)\\ Embedding personal-data protection within a dedicated internal function, the Biocodex, a global pharmaceutical company focused on probiotics and specialty healthcare, mandate demonstrates the team’s role in translating data-protection requirements into an organised institutional structure rather than treating them as a standalone formal exercise.

Biocodex instructed the team to advise on the design and implementation of a structured personal data protection framework. The work centred on establishing the foundation for a coordinated approach to personal-data protection, with the framework serving as the underlying structure for the client’s internal treatment of these requirements. The mandate therefore combined the design phase with implementation, ensuring that the proposed structure was carried into the client’s internal organisation.

Arturo González de Araujo, Ernesto Zavala and Damián Ortega led the engagement. Their contribution focused on advising Biocodex through the creation and implementation of the structured personal data protection framework and on aligning that work with the establishment of a dedicated Personal Data Protection Department.

The engagement resulted in the establishment of a dedicated Personal Data Protection Department, a concrete institutional outcome that strengthened internal controls and supported the reduction of regulatory exposure. The matter provides evidence of the team’s ability to help a client move from a general personal-data protection requirement to a dedicated department and a structured framework for its implementation.
\item \textbf{Hotel Riazor Data Protection Enhancement} (Client: Hotel Riazor | Value: N/A)\\ Representing personal data protection within a dedicated internal function gives this mandate significance beyond the adoption of isolated policies: Hotel Riazor, a business hotel offering lodging and event services in Mexico City, sought to reinforce the institutional arrangements through which lawful data handling is managed and governed. The engagement centred on strengthening Hotel Riazor’s personal data protection and governance framework, with the work directed at the client’s internal approach to the treatment of personal data.

The firm supported Hotel Riazor in reinforcing its Personal Data Department, recognising that the department’s role was central to the client’s ability to administer personal-data responsibilities through an identifiable internal structure. Alongside that departmental reinforcement, the team implemented internal policies designed to ensure lawful data handling. The mandate therefore combined organisational work within the Personal Data Department with the implementation of documented internal policies, linking governance arrangements to the day-to-day handling of personal data.

Arturo González de Araujo, Ernesto Zavala and Damián Ortega led the work. Their contribution encompassed the advice on strengthening Hotel Riazor’s personal data protection and governance framework, the support provided to reinforce the Personal Data Department, and the implementation of internal policies for lawful data handling. The resulting institutional measures were the reinforcement of Hotel Riazor’s Personal Data Department and the implementation of internal policies intended to ensure that personal data is handled lawfully.
\item \textbf{Grupo Excelsior Data Protection Regularisation} (Client: Grupo Excelsior | Value: N/A)\\ Regularisation and restructuring of operations can make data protection an operational priority rather than a discrete documentation exercise; for Grupo Excelsior, one of Mexico’s leading dairy producers, data protection regularisation is presented as such a mandate. The team advised Grupo Excelsior on the regularisation and restructuring of its operations, with data protection as the specific focus of the engagement. This positioned the work at the point where operational arrangements required formal reorganisation and where data-protection requirements needed to be incorporated into the client’s internal approach.

Arturo González de Araujo, Ernesto Zavala and Damián Ortega led the work for Grupo Excelsior. Their specific role extended beyond identifying the data-protection focus of the restructuring: the firm designed and implemented internal policies and training programmes. The mandate therefore combined the regularisation of operations, the restructuring of those operations, the creation and implementation of internal policies, and training programmes intended to support the application of the new internal arrangements.

The stated outcome was the design and implementation of internal policies and training programmes in connection with the regularisation and restructuring of operations focusing on data protection. For Grupo Excelsior, these measures, significantly strengthening compliance posture and reducing regulatory exposure, constituted the concrete result of the engagement. The matter evidences a data-protection regularisation exercise in which Arturo González de Araujo, Ernesto Zavala and Damián Ortega translated the client’s operational restructuring requirements into implemented internal policies and training programmes.
\item \textbf{Grupo Modelquipo Data Protection Framework} (Client: Grupo Modelquipo | Value: N/A)\\ Assisting Grupo Modelquipo, industrial equipment and machinery solutions provider, Data Protection Framework places data protection within the client’s corporate governance strategy, evidencing an implementation-focused mandate in which privacy considerations were connected directly to risk management, compliance and governance culture rather than addressed as an isolated exercise.

The submission describes the client context as: Grupo Modelquipo, improving regulatory compliance and integrating data protection into governance culture. Grupo Modelquipo instructed the firm to advise on the implementation of a data protection framework as part of its corporate governance strategy. The firm’s specific role was therefore to support implementation of that framework while focusing the work on risk management and compliance, and on the incorporation of data protection into the client’s governance culture.

Arturo González de Araujo, Ernesto Zavala and Damián Ortega led the mandate. Their work connected the implementation of the data protection framework to Grupo Modelquipo’s corporate governance strategy, with risk management and compliance forming the central focus of the engagement. The stated result was improved regulatory compliance and the integration of data protection into governance culture, achieved through the implementation of the data protection framework as part of the client’s corporate governance strategy.
\item \textbf{Tiendas Chedraui Data Protection Advisory} (Client: Tiendas Chedraui | Value: N/A)\\ At the point where new store openings and operational expansion are translated into operational decisions, Tiendas Chedraui, a Mexican supermarket and retail chain, has sought partner-led legal assessment of the data-processing risks associated with those developments. The significance of the mandate lies in placing the evaluation of data-processing risk alongside proposed growth activity, so that the implications of new openings and broader operational expansion are considered as part of the relevant legal review.

Arturo González de Araujo and Damián Ortega advise Tiendas Chedraui on the legal assessment of new store openings and operational expansion. Their specific role is to evaluate the risks associated with data processing arising in those operational contexts, connecting the client’s proposed developments with the regulatory requirements and governance standards that apply to the relevant processing activities. The work is directed not merely at identifying the legal considerations linked to expansion, but at assessing those considerations in relation to the client’s proposed store-opening and operational decisions.

For Tiendas Chedraui, ensuring alignment with regulatory requirements and governance standards is the concrete result pursued through the assessment of data-processing risks connected with each new store opening and with operational expansion. The mandate therefore evidences a defined legal role for Arturo González de Araujo and Damián Ortega: evaluating data-processing risk in the context of the client’s new openings and expansion, and ensuring alignment between those activities, regulatory requirements and governance standards.

\end{itemize}

\section*{IV. Strategic Evolution Path}
\begin{itemize}
\item Maintain current strategy.
\end{itemize}

\vfill
\centering
\small{This document is a structural assessment and does not guarantee ranking results.}

\end{document}